\section{Signal and Execution Model}
\label{sec:signal-model}
Let $x[n]$ be a real signal sampled at $f_s$ Hz. A sample call advances one
engine sample; a host callback can render $D$ such samples. These are distinct
measurement boundaries. A hop $H$ specifies the interval between selected
frame endpoints, not the host callback size. The transform length $N\ge4$ is
a power of two, and $H\ge1$ is an integer. With continuously captured input
and fixed hop, the frame endpoint is $t_r=rH$. The selected window is
\begin{equation}
 u_r[m]=x[t_r-N+1+m]w[m],\qquad 0\le m<N,
\end{equation}
where $w$ is an analysis window. Its unnormalized forward transform is
\begin{equation}
 X_r[k]=\sum_{m=0}^{N-1}u_r[m]e^{-j2\pi km/N}.
 \label{eq:dft}
\end{equation}
The associated bin frequency is $kf_s/N$. Deferring computation changes when
$X_r$ becomes available; it does not change which samples enter this sum.
Sample-index age measures the distance from a specified frame timestamp to
publication, converted to seconds using $f_s$. It excludes wall-clock
computation, device buffering, and display consumption. In particular, a
transform computed in its endpoint's sample call has zero endpoint-to-publication
age but nonzero execution time.

For amplitude calibration, the finite-window coherent gain is
$G_c=N^{-1}\sum_m w[m]$. An isolated interior bin-centered real sinusoid can be
scaled using $2|X_r[k]|/(NG_c)$ when the window's positive- and negative-frequency
contributions do not overlap materially. DC and Nyquist do not receive the
factor of two. This is amplitude calibration, not power-spectral-density
normalization; the latter involves window energy. The plugin caches window
samples and applies tabulated coherent-gain factors. Those factors should not
be assumed to equal the exact finite-length sum for every window and length.
The original FFT campaign uses explicit rectangular and periodic Hann samples
and does not certify display calibration. The pipeline campaign instead uses
a coherent-gain-normalized periodic Hann window.

Real-input packing uses $M=N/2$ complex samples;
the analyzer retains $K=M+1$ nonnegative-frequency bins, including DC and
Nyquist. The standard radix-2 traversal requires
$B=B_r(N)=(N/4)\log_2(N/2)$ butterflies. Appendix~\ref{sec:legacy-transform}
preserves the butterfly equations, arithmetic-order proof, resumable state
transition, and real-spectrum reconstruction.
